%% =========================================================================
%% DRAFT replacement for the material from line ~5736 of monoidal.tex
%% (everything from the definition of \llangle-\rrangle_{\join} down to and
%% including the three half-written theorems labelled
%% theorem:edpoi_rewriting_frob).
%%
%% Architecture:
%%   \mathbf{H}_\Sigma[\join]  --E-->   EHypI_*(\Sigma)
%%          |                              |
%%          | (SLat quotient)              | Q
%%          v                              v
%%   \mathbf{H}^\join_\Sigma  --<<->>-- EHypI_*(\Sigma)/R
%%
%% E is the concrete, term-level counterpart of Bonchi's \llangle-\rrangle;
%% \llangle-\rrangle_{\join} is recovered as the induced map on quotients.
%% =========================================================================

\subsection{Raw \texorpdfstring{$\join$}{join}-expressions}

The interpretation $\llangle - \rrangle_{\join}$ of~\cref{theorem:ehypi_slat_isomorphism}
takes values in the \emph{quotient} $\catname{EHypI}_{\bullet}(\Sigma)/\mathfrak{R}$, so
$\llangle \corner{a} \rrangle_{\join}$ is a $\sim_{\mathfrak{R}}$-class of extended
cospans rather than a single one.
EDPOI rewriting (\cref{def:dpoi_rooted}), by contrast, is defined on individual
well-typed extended cospans, and it is individual extended cospans---not their
classes---that a machine holds in memory.
This is a genuine departure from the situation of~\cref{chapter:string_diagrams}.
There, $\llangle - \rrangle : \mathbf{H}_{\Sigma} \to \catname{HypI}(\Sigma)$ is an
isomorphism onto a category of concrete objects, because the equations of a symmetric
monoidal category and of a Frobenius algebra hold \emph{on the nose} in
$\catname{HypI}(\Sigma)$.
In $\catname{EHypI}_{\bullet}(\Sigma)$ the corresponding statement is only partially
true: the symmetric monoidal and Frobenius equations are absorbed by the
representation, and so is commutativity of $\join$ (\cref{remark:commutativity_absorbed}),
but associativity, idempotency, distributivity and reindexing are not---they are
precisely the content of the structural system $\mathfrak{R}$
(\cref{def:structural_rules}).

We therefore separate the two roles that $\mathbf{H}^{\join}_{\Sigma}$ has been playing.
The syntax of $\catname{SLat}$-string diagrams---the object one draws, and the object
one represents---is a category of \emph{raw} expressions, in which $\join$ satisfies no
equations; the $\catname{SLat}$-PROP $\mathbf{H}^{\join}_{\Sigma}$ is its quotient.
The interpretation of raw expressions is concrete, and the interpretation
$\llangle - \rrangle_{\join}$ is what it induces on quotients.

The ambient notion is the following.
It is deliberately \emph{weaker} than $\catname{SLat}$-enrichment: the join is extra
structure carried by the hom-sets, and neither $\fatsemi$ nor $\otimes$ is required to
respect it.

\begin{definition}[$\catname{PROP}[\join]$]\label{def:prop_with_raw_join}
Write $M_{\geq 2}$ for the endofunctor on $\catname{Set}$ sending a set $X$ to the set of
finite multisets of elements of $X$ of cardinality at least $2$.
A \emph{$\catname{PROP}$ with raw join} is a PROP $\mathbb{C}$ together with, for each
pair of objects $n, m$, a function
\[
\join_{n,m} : M_{\geq 2}\bigl(\mathbb{C}(n,m)\bigr) \longrightarrow \mathbb{C}(n,m)~,
\]
subject to no equations whatsoever.
A morphism of such is an identity-on-objects strict symmetric monoidal functor commuting
with the $\join_{n,m}$.
We write $\catname{PROP}[\join]$ for the resulting category.
\end{definition}

\begin{remark}\label{remark:raw_join_vs_slat}
Three differences from~\cref{def:slat-prop} should be kept in view.
First, the hom-objects are algebras for the \emph{endofunctor} $M_{\geq 2}$, not for a
monad: there is no associativity, so a nested join is not a join.
Second, and most importantly, an $\catname{SLat}$-PROP requires composition and tensor to
be $\catname{SLat}$-morphisms in each variable---that is, to distribute over
$\join$---whereas~\cref{def:prop_with_raw_join} imposes no such condition; distributivity
is one of the rules of $\mathfrak{R}$.
A $\catname{PROP}$ with raw join is therefore not enriched in anything.
Third, the forgetful functor $\catname{PROP}[\join] \to \catname{PROP}$ has a left adjoint
$\mathrm{Free}_{\join}$, and \cref{def:raw_join_expressions} below is
$\mathrm{Free}_{\join}(\mathbf{H}_{\Sigma})$; this adjunction is the raw-level counterpart
of the free-enrichment adjunction of~\cref{prop:prop-adjunction}, and the two are related
by the quotient $\pi$ constructed after~\cref{def:raw_join_expressions}.
\end{remark}

\begin{definition}[Raw $\join$-expressions]\label{def:raw_join_expressions}
Let $\mathbf{H}_{\Sigma}[\join]$ be the strict symmetric monoidal category with objects
the natural numbers, and with morphisms $n \to m$ generated by the grammar
\[
A, B \;::=\; \eta(h) \;\mid\; A \fatsemi B \;\mid\; A \otimes B \;\mid\;
\bigjoin\{\!\{A_{1}, \ldots, A_{p}\}\!\}\quad (p \geq 2)~,
\]
where $h$ ranges over the morphisms of $\mathbf{H}_{\Sigma}$, the typing rules are the
evident ones, all $A_{r}$ in a join have the same domain and codomain as the join
itself, and $\{\!\{-\}\!\}$ denotes a finite \emph{multiset}, so that the operands of a
join carry no order but may repeat.
Morphisms are taken modulo the equations of a strict symmetric monoidal category
together with the requirement that $\eta$ be a strict symmetric monoidal functor, i.e.
\[
\eta(h_{1}) \fatsemi \eta(h_{2}) = \eta(h_{1} \fatsemi h_{2})~, \qquad
\eta(h_{1}) \otimes \eta(h_{2}) = \eta(h_{1} \otimes h_{2})~, \qquad
\eta(\id_{n}) = \id_{n}~, \qquad \eta(\sym_{n,m}) = \sym_{n,m}~.
\]
No equation relates $\join$ to $\fatsemi$, to $\otimes$, or to itself.
We call the morphisms of $\mathbf{H}_{\Sigma}[\join]$ \emph{raw
$\join$-expressions}, and write $\mathbf{S}_{\Sigma}[\join]$ for the subcategory
generated by $\eta(\mathbf{S}_{\Sigma})$.
\end{definition}

\begin{remark}[Why the join is $p$-ary, unordered and repeating]\label{remark:join_arity}
The three choices are each forced by the representation, and it is worth recording why a
binary commutative join---the first thing one would try---does not work.

\emph{Arity.}
By non-degeneracy (\cref{def:rooted_ehypergraph}) a $\boxsort$-edge has at least two
component children, but nothing bounds their number, so a box with three components is a
legitimate carrier.
A binary $\join$ produces, by~\cref{def:join_of_extended_cospans}, only two-component
boxes, and $\mathcal{E}$ would fail to be surjective.
Nor can this be repaired by imposing associativity on raw expressions:
$\mathcal{F} \join (\mathcal{G} \join \mathcal{H})$ builds a box one of whose components is
the root $\top_{\mathcal{G} \join \mathcal{H}}$ and therefore itself contains a box, which
is not isomorphic to the flat three-component box built by the ternary join.
Associativity is thus a rule of $\mathfrak{R}$---and, by the same token, its instances in
\cref{def:structural_rules} must be read as \emph{flattening} rules at every arity
$p \geq 2$, merging a nested box into its parent, since it is flat joins that
\cref{lemma:canonical_representative} produces as canonical representatives.

\emph{Repetition.}
Operands form a multiset rather than a set because idempotency is likewise a rule of
$\mathfrak{R}$: $A \join A$ denotes a box with two isomorphic components, which is not
isomorphic to $\mathcal{E}(A)$.

\emph{Order.}
Operands are unordered because commutativity \emph{is} absorbed by the representation
(\cref{remark:commutativity_absorbed}), the components of a $\boxsort$-edge carrying no
order.
\end{remark}

There is an evident identity-on-objects strict symmetric monoidal functor
\[
\pi : \mathbf{H}_{\Sigma}[\join] \longrightarrow
\widetilde{U}_{\catname{SLat}}(\mathbf{H}^{\join}_{\Sigma})~,
\]
sending $\eta(h)$ to $\{h\}$ and $\bigjoin\{\!\{A_{r}\}\!\}$ to
$\bigjoin_{r} \pi(A_{r})$; it is surjective on morphisms
by~\cref{rem:slat-free-concrete}.
We write $A =_{\join} B$ for $\pi(A) = \pi(B)$.

\subsection{The concrete interpretation}

\begin{definition}[Interpretation of raw $\join$-expressions]\label{def:concrete_interpretation}
Define
\[
\mathcal{E} : \mathbf{H}_{\Sigma}[\join] \longrightarrow \catname{EHypI}_{\bullet}(\Sigma)
\]
to be the identity on objects and, on morphisms, by structural recursion:
\begin{align*}
\mathcal{E}(\eta(h)) &\defeq \Xi\bigl(\llangle h \rrangle\bigr)~,\\
\mathcal{E}(A \fatsemi B) &\defeq \mathcal{E}(A) \fatsemi \mathcal{E}(B)~,\\
\mathcal{E}(A \otimes B) &\defeq \mathcal{E}(A) \otimes \mathcal{E}(B)~,\\
\mathcal{E}\bigl(\bigjoin\{\!\{A_{1}, \ldots, A_{p}\}\!\}\bigr)
  &\defeq \mathcal{E}(A_{1}) \join \cdots \join \mathcal{E}(A_{p})~,
\end{align*}
where $\Xi$ is the embedding of~\cref{prop:embedding_hypi} and the last clause is the
$p$-ary join of~\cref{def:join_of_extended_cospans}.
\end{definition}

That $\mathcal{E}$ is well defined on the equations of~\cref{def:raw_join_expressions} is
immediate: $\Xi \circ \llangle - \rrangle$ is a strict symmetric monoidal functor by
\cref{prop:embedding_hypi} and~\cref{theorem:frobenius_isomorphism}, and the $p$-ary join
is invariant under permutation of its operands
by~\cref{def:isomorphic_ecospans,remark:commutativity_absorbed}.

\begin{theorem}[Concrete representation]\label{theorem:concrete_representation}
$\mathcal{E}$ is an isomorphism of symmetric monoidal categories
\[
\mathbf{H}_{\Sigma}[\join] \;\cong\; \catname{EHypI}_{\bullet}(\Sigma)
\]
which moreover commutes with $\join$, i.e.\
$\mathcal{E}(\bigjoin\{\!\{A_{r}\}\!\}) = \bigjoin_{r} \mathcal{E}(A_{r})$.
\end{theorem}

The proof is not a direct appeal to~\cref{theorem:frobenius_isomorphism}: that theorem
concerns PROPs, and a $\boxsort$-edge is not a generator of $\Sigma$.
It is instead an induction on hierarchical depth (\cref{def:depth}), in which
\cref{theorem:frobenius_isomorphism} is applied at each level to an \emph{enlarged}
signature, obtained by treating the boxes occurring at that level as generators.

\begin{definition}[Box signature]\label{def:box_signature}
For a PROP $\mathbb{X}$, let $\mathrm{Box}(\mathbb{X})$ be the monoidal signature with, for
each $n, m$, one generator $n \to m$ for every finite multiset of cardinality at least $2$
of elements of $\mathbb{X}(n,m)$.
For a set $B$ of $\boxsort$-edges of a rooted e-hypergraph we write $\Sigma_{B}$ for the
signature $\Sigma$ enlarged by one generator $|s(b)| \to |t(b)|$ for each $b \in B$.
\end{definition}

\begin{proof}[Proof of~\cref{theorem:concrete_representation}]
Both categories are identity-on-objects over $\mathbb{N}$, so it suffices to show that
each hom-map is a bijection.
We construct the inverse by induction on the hierarchical depth of the carrier, which is
finite because carriers are finite (\cref{def:rooted_ehypergraph}).

Let
\[
n \xrightarrow{f_{\mathrm{ext}}} n' \xrightarrow{f_{\mathrm{int}}} \mathcal{G}
  \xleftarrow{g_{\mathrm{int}}} m' \xleftarrow{g_{\mathrm{ext}}} m
\]
be a well-typed extended cospan and let $B$ be the set of $\boxsort$-edges $b$ of
$\mathcal{G}$ with $\top_{\mathcal{G}} <^{\mu}_{\mathcal{G}} b$, that is, the boxes at the
top level.
Deleting from $\mathcal{G}$ every element strictly below some $b \in B$, and then
forgetting $\top_{\mathcal{G}}$ and $<^{\mu}$, yields an ordinary hypergraph
$G^{\flat}$ over $\Sigma_{B}$.
By the last condition of~\cref{def:extended_cospan} every strictly internal interface
vertex is deleted in this process, so the external legs make
\[
n \xrightarrow{(f_{\mathrm{ext}} \fatsemi f_{\mathrm{int}})_{V}} G^{\flat}
  \xleftarrow{(g_{\mathrm{ext}} \fatsemi g_{\mathrm{int}})_{V}} m
\]
a morphism of $\catname{HypI}(\Sigma_{B})$.
By~\cref{theorem:frobenius_isomorphism}, applied to the signature $\Sigma_{B}$, it is
$\llangle t \rrangle$ for a unique $t \in \mathbf{H}_{\Sigma_{B}}(n,m)$.

Let $b \in B$ have components $c_{1}, \ldots, c_{p}$, where $p \geq 2$ by
non-degeneracy.
The sub-e-hypergraph rooted at $c_{r}$ (\cref{def:sub_ehypergraph}), equipped with the
blocks $\pi_{f}^{-1}(c_{r})$ and $\pi_{g}^{-1}(c_{r})$ of~\cref{def:isomorphic_ecospans}
as external interfaces and with the blocks indexed by $\cmpsort$-edges strictly below
$c_{r}$ as strictly internal interfaces, is a well-typed extended cospan
$\mathcal{G}_{b,r} : |s(b)| \to |t(b)|$: well-typedness supplies exactly the bijections
between those blocks and the sequences $s(b)$ and $t(b)$ discussed
after~\cref{def:isomorphic_ecospans}.
Its depth is strictly smaller than that of $\mathcal{G}$, so by the induction hypothesis
$\mathcal{G}_{b,r} = \mathcal{E}(A_{b,r})$ for a unique raw $A_{b,r}$.

Substituting $\bigjoin\{\!\{A_{b,1},\ldots,A_{b,p}\}\!\}$ for the generator $b$ in $t$,
for each $b \in B$, produces a raw expression $A$ with $\mathcal{E}(A) \cong \mathcal{G}$.
This substitution inverts the clauses of~\cref{def:concrete_interpretation}: the
$\join$-clause because~\cref{def:join_of_extended_cospans} builds a box whose top level
consists of the single edge $b$ together with its wires, and the $\fatsemi$- and
$\otimes$-clauses by~\cref{lemma:boxes_preserved} below.

Injectivity follows by the same induction.
The term $t$ is determined by $G^{\flat}$ up to the equations of
$\mathbf{H}_{\Sigma_{B}}$, which are exactly those imposed on raw expressions
in~\cref{def:raw_join_expressions}; and the multiset $\{\!\{A_{b,r}\}\!\}_{r}$ is
determined by the components of $b$, which carry no order.
Functoriality and strict monoidality hold by construction.
\end{proof}

The following is the one point at which the induction above needs the specific shape of
composition in $\catname{EHypI}_{\bullet}(\Sigma)$, and it deserves to be proved rather
than assumed.

\begin{lemma}\label{lemma:boxes_preserved}
Let $\mathcal{F}$ and $\mathcal{G}$ be well-typed extended cospans.
The top-level $\boxsort$-edges of $\mathcal{F} \fatsemi \mathcal{G}$ are the disjoint union
of those of $\mathcal{F}$ and of $\mathcal{G}$, and each retains its multiset of
components; likewise for $\mathcal{F} \otimes \mathcal{G}$.
Consequently the passage $\mathcal{G} \mapsto G^{\flat}$ of the previous proof carries
$\fatsemi$ and $\otimes$ of extended cospans to $\fatsemi$ and $\otimes$ in
$\catname{HypI}(\Sigma_{B})$.
\end{lemma}
\begin{proof}[Proof sketch]
Composition is the pushout of~\cref{prop:pushout_exists} along a span whose apex is
pointed discrete, so the relations $R_{V}$ and $R_{E}$ of that construction identify only
vertices and the two roots; in particular no two distinct $\boxsort$-edges are identified,
no $\boxsort$-edge is identified with a $\cmpsort$-edge, and the immediate children of a
$\boxsort$-edge are untouched.
The claim for $\otimes$ is immediate, the monoidal product being computed by the same
pushout over $\mathfrak{T}$.
\end{proof}

\begin{remark}[Fixed-point reading]\label{remark:fixed_point}
The induction above is an unfolding of a recursive equation: writing
$\Theta(\mathbb{X}) \defeq \catname{HypI}(\Sigma + \mathrm{Box}(\mathbb{X}))$, the proof
exhibits mutually inverse maps
\[
\catname{EHypI}_{\bullet}(\Sigma)
\;\cong\;
\Theta\bigl(\catname{EHypI}_{\bullet}(\Sigma)\bigr)~,
\qquad
\mathbf{H}_{\Sigma}[\join] \;\cong\; \Theta\bigl(\mathbf{H}_{\Sigma}[\join]\bigr)~,
\]
and $\mathcal{E}$ is the induced comparison.
The induction terminates on the \emph{least} such fixed point because carriers are finite
and raw expressions are finite terms; permitting infinite nesting on either side would
break the argument, and one would be comparing greatest fixed points instead.
This reading is not needed below, but it explains why the recursion is structural rather
than ad hoc, and why finiteness in~\cref{def:rooted_ehypergraph} is doing real work.
\end{remark}

\begin{remark}
\Cref{theorem:concrete_representation} is the exact counterpart
of~\cref{theorem:frobenius_isomorphism}: an isomorphism between a category of
\emph{syntax} and a category of \emph{concrete combinatorial objects}, with the
equations of the ambient structure absorbed on the nose.
What is absorbed here is strictly less than in the Frobenius case, and the deficit
is measured by $\mathfrak{R}$.
\Cref{example:non_isomorphic_cospans} and the example following it are instances: the
raw expressions $\sym \fatsemi (\sym \join (\sym \fatsemi (f \otimes g)))$ and
$(\id \otimes \id) \join (f \otimes g)$ are distinct in
$\mathbf{H}_{\Sigma}[\join]$ and interpreted by non-isomorphic extended cospans, while
being equal in $\mathbf{H}^{\join}_{\Sigma}$.

The induction of~\cref{theorem:concrete_representation} also isolates what changes in
the monogamous acyclic setting: replacing $\catname{HypI}$ by $\catname{MHypI}$ throughout,
and~\cref{theorem:frobenius_isomorphism} by~\cref{theorem:monogamous_isomorphism}, gives
the corresponding statement for $\mathbf{S}_{\Sigma}[\join]$, once the monogamy and
convexity conditions have been formulated for extended cospans.
\end{remark}

\begin{proposition}[Factorisation]\label{prop:interpretation_factorisation}
The square
\[\begin{tikzcd}
	{\mathbf{H}_{\Sigma}[\join]} && {\catname{EHypI}_{\bullet}(\Sigma)} \\
	\\
	{\mathbf{H}^{\join}_{\Sigma}} && {\catname{EHypI}_{\bullet}(\Sigma)/\mathfrak{R}}
	\arrow["{\mathcal{E}}", from=1-1, to=1-3]
	\arrow["\pi"', from=1-1, to=3-1]
	\arrow["Q", from=1-3, to=3-3]
	\arrow["{\llangle - \rrangle_{\join}}"', from=3-1, to=3-3]
\end{tikzcd}\]
commutes, and $\mathcal{E}$ both preserves and reflects the two quotients:
for all raw $A, B$,
\[
A =_{\join} B
\qquad\Longleftrightarrow\qquad
\mathcal{E}(A) \sim_{\mathfrak{R}} \mathcal{E}(B)~.
\]
\end{proposition}
\begin{proof}[Proof sketch]
Commutativity is checked on the generators of~\cref{def:raw_join_expressions}, using
that $\llangle - \rrangle_{\join}$ was defined as
$\widetilde{F}_{\catname{SLat}}(\llangle - \rrangle) \fatsemi \Phi$ and that
$\Phi(\{F\}) = [\Xi(F)]$.
The forward implication of the displayed equivalence is soundness of $\mathfrak{R}$ for
the $\catname{SLat}$ equations: each equation of~\cref{fig:string-equations} other than
commutativity is an instance of a rule of $\mathfrak{R}$ by
\cref{def:structural_rules}, commutativity is absorbed, and $\sim_{\mathfrak{R}}$ is a
congruence for $\fatsemi$, $\otimes$ and $\join$
by~\cref{lemma:rewriting_congruence}.
The backward implication is completeness of $\mathfrak{R}$, which is
\cref{lemma:canonical_representative}: $\mathcal{E}(A) \sim_{\mathfrak{R}} \mathcal{E}(B)$
forces the two canonical representatives to coincide, hence
$\operatorname{Flat}(\mathcal{E}(A)) = \operatorname{Flat}(\mathcal{E}(B))$, and by
\cref{prop:embedding_hypi} this set is exactly $\pi(A) = \pi(B)$ transported along
$\llangle - \rrangle$.
\end{proof}

\begin{corollary}
\Cref{theorem:ehypi_slat_isomorphism} follows: $\llangle - \rrangle_{\join}$ is the
isomorphism induced on quotients by the isomorphism $\mathcal{E}$.
\end{corollary}

\subsection{Rewriting}

We first lift $\Rightarrow$ of~\cref{def:enriched_prop_rewriting} to raw expressions.

\begin{definition}[Raw contexts and raw rewriting]\label{def:raw_rewriting}
A \emph{raw context} $C[-] : (i \to j) \Rightarrow (m \to n)$ is a raw expression with a
single hole, generated by
\[
C[-] \;::=\; [-] \;\mid\; A \fatsemi C[-] \;\mid\; C[-] \fatsemi A \;\mid\;
A \otimes C[-] \;\mid\; C[-] \otimes A \;\mid\;
\bigjoin\{\!\{C[-], A_{2}, \ldots, A_{p}\}\!\}~.
\]
Given raw $L, R : i \to j$, we write $A \Rightarrow_{\langle L, R \rangle} B$ if
$A = C[L]$ and $B = C[R]$ for some raw context $C[-]$.
\end{definition}

\begin{remark}
The two clauses of~\cref{def:enriched_prop_rewriting} are exactly the two shapes of raw
context: the first, $a = a_{1} \fatsemi (\id_{k} \otimes l) \fatsemi a_{2}$, is the
general context built from $\fatsemi$ and $\otimes$ by the equations of a symmetric
monoidal category, and the second is the $\join$-clause, whose hole sits in one operand of
a join.
Thus $\pi$ carries $\Rightarrow$ of~\cref{def:raw_rewriting} onto $\Rightarrow$ of
\cref{def:enriched_prop_rewriting}.
\end{remark}

The following is the $\catname{SLat}$ counterpart
of~\cref{theorem:dpoi_rewriting_frob}, and it involves no quotient on either side.

\begin{theorem}[Soundness and completeness, raw form]\label{theorem:edpoi_raw}
Let $L, R : i \to j$ and $A, B : m \to n$ be raw $\join$-expressions.
Then
\[
A \Rightarrow_{\langle L, R \rangle} B
\qquad\Longleftrightarrow\qquad
\mathcal{E}(\corner{A})
\;\overline{\rightsquigarrow}_{\langle \mathcal{E}(\corner{L}), \mathcal{E}(\corner{R}) \rangle}\;
\mathcal{E}(\corner{B})~,
\]
where $\overline{\rightsquigarrow}$ is EDPOI rewriting of~\cref{def:dpoi_rooted}.
\end{theorem}

\begin{proof}[Proof strategy]
Both directions go by induction on the number of $\join$-clauses in the context, resp.\
on the hierarchical depth of the anchor of the match, the two being matched by
\cref{lemma:edpoi_locality}.

\emph{Only if.}
If $C[-]$ has no $\join$-clause then $A = A_{1} \fatsemi (\id_{k} \otimes L) \fatsemi A_{2}$
and, applying $\corner{-}$ and $\mathcal{E}$, the argument
of~\cref{theorem:dpoi_rewriting_frob} applies verbatim: composition in
$\catname{EHypI}_{\bullet}(\Sigma)$ is given by the pushout
of~\cref{prop:pushout_exists} (\cref{def:rooted_ehypi}), so the folded decomposition
$\mathcal{E}(\corner{A}) = \mathcal{E}(\corner{L}) \fatsemi \mathcal{E}(\widetilde{A})$
exhibits the two EDPOI squares, with $\mathcal{E}(\widetilde{A})$ the extended pushout
complement.
Conditions~(1)--(5) of~\cref{def:extended_complement} hold because
$\mathcal{E}(\widetilde{A})$ is an extended cospan and its interface legs are strict, so
the match is anchored at $\top$.
If the outermost clause of $C[-]$ is
$\bigjoin\{\!\{C'[-], A_{2}, \ldots, A_{p}\}\!\}$, apply the induction hypothesis to
$C'[L] \Rightarrow C'[R]$ and then clause~(3)
of~\cref{lemma:rewriting_congruence}; the remaining clauses of $C[-]$ are handled by
clauses~(1) and~(2) of the same lemma.

\emph{If.}
Let $c$ be the anchor of the match $m : \mathcal{E}(\corner{L}) \to
\mathcal{E}(\corner{A})$.
If $c = \top$, \cref{lemma:edpoi_locality} folds the extended pushout complement into a
well-typed extended cospan $\mathcal{C}$ with
$\mathcal{E}(\corner{A}) = \mathcal{E}(\corner{L}) \fatsemi \mathcal{C}$ and
$\mathcal{E}(\corner{B}) = \mathcal{E}(\corner{R}) \fatsemi \mathcal{C}$
\emph{on the nose}, not merely up to $\mathfrak{R}$.
By~\cref{theorem:concrete_representation} there is a unique raw $\widetilde{A}$ with
$\mathcal{E}(\widetilde{A}) = \mathcal{C}$, and unbending the wires as
in~\cref{theorem:dpoi_rewriting_frob} yields the required $\fatsemi$/$\otimes$ context.
If $c \neq \top$, let $b = {<^{\mu}}(c)$.
By~\cref{lemma:edpoi_locality} the box $b$ and all its components other than $c$ are
untouched, and $\mathcal{E}(\corner{A})[b \mapsto c] \;\overline{\rightsquigarrow}\;
\mathcal{E}(\corner{B})[b \mapsto c]$ has an anchor of strictly smaller depth; the
induction hypothesis supplies a raw context for the resolved rewrite, and reinstating $b$
wraps it in a $\join$-clause whose remaining operands are the raw expressions
corresponding, by~\cref{theorem:concrete_representation}, to the untouched components.
\end{proof}

\begin{remark}[Why no modulo appears]\label{remark:no_modulo_raw}
\Cref{theorem:edpoi_raw} is sharper than a statement about $\sim_{\mathfrak{R}}$-classes
would be, and it is the statement one wants for an implementation: it says that a single
EDPOI step on the extended cospan one actually holds corresponds to a single rule
application on the diagram it denotes, with no structural steps interposed.
The structural rules of $\mathfrak{R}$ are thereby demoted from part of the semantics to
operations one may \emph{choose} to perform---flattening, sharing, congruence closure---and
in particular one is never obliged to normalise, which matters because the canonical form
of~\cref{lemma:canonical_representative} is the maximally unshared one.
\end{remark}

Passing to $\catname{SLat}$-morphisms reintroduces the quotient on both sides, and only
there.

\begin{definition}[Rewriting modulo $\mathfrak{R}$]\label{def:rewriting_modulo}
For well-typed extended cospans $\mathcal{X}, \mathcal{Y}$ and a rewrite rule
$\rho = \langle \mathcal{L}, \mathcal{R} \rangle$, write
\[
\mathcal{X} \;\overline{\rightsquigarrow}^{\,\mathfrak{R}}_{\rho}\; \mathcal{Y}
\]
if there exist $\mathcal{X}', \mathcal{Y}'$ with
$\mathcal{X} \sim_{\mathfrak{R}} \mathcal{X}'
 \;\overline{\rightsquigarrow}_{\rho}\;
 \mathcal{Y}' \sim_{\mathfrak{R}} \mathcal{Y}$.
\end{definition}

\begin{theorem}[Soundness and completeness]\label{theorem:edpoi_rewriting_slat}
Let $\langle l, r \rangle$ be a rewrite rule in $\mathbf{H}^{\join}_{\Sigma}$ and let
$a, b : m \to n$ in $\mathbf{H}^{\join}_{\Sigma}$.
Then the following are equivalent.
\begin{enumerate}
  \item $a \Rightarrow_{\langle l, r \rangle} b$ in the sense
        of~\cref{def:enriched_prop_rewriting};
  \item there exist raw $A, B, L, R$ with $\pi(A) = a$, $\pi(B) = b$, $\pi(L) = l$,
        $\pi(R) = r$ and
        \[
        \mathcal{E}(\corner{A})
        \;\overline{\rightsquigarrow}^{\,\mathfrak{R}}_{\langle \mathcal{E}(\corner{L}), \mathcal{E}(\corner{R}) \rangle}\;
        \mathcal{E}(\corner{B})~;
        \]
  \item for \emph{every} raw $A$ with $\pi(A) = a$ and every raw $L, R$ with
        $\pi(L) = l$, $\pi(R) = r$, there is a raw $B$ with $\pi(B) = b$ and
        \[
        \mathcal{E}(\corner{A})
        \;\overline{\rightsquigarrow}^{\,\mathfrak{R}}_{\langle \mathcal{E}(\corner{L}), \mathcal{E}(\corner{R}) \rangle}\;
        \mathcal{E}(\corner{B})~.
        \]
\end{enumerate}
\end{theorem}
\begin{proof}
$(3) \Rightarrow (2)$ is immediate, since $\pi$ is surjective on morphisms.

$(2) \Rightarrow (1)$.
Suppose $\mathcal{E}(\corner{A}) \sim_{\mathfrak{R}} \mathcal{X}'
\;\overline{\rightsquigarrow}\; \mathcal{Y}' \sim_{\mathfrak{R}} \mathcal{E}(\corner{B})$.
By~\cref{theorem:concrete_representation} write $\mathcal{X}' = \mathcal{E}(\corner{A'})$
and $\mathcal{Y}' = \mathcal{E}(\corner{B'})$; by
\cref{prop:interpretation_factorisation}, $\pi(A') = \pi(A) = a$ and
$\pi(B') = \pi(B) = b$.
By~\cref{theorem:edpoi_raw}, $A' \Rightarrow_{\langle L, R \rangle} B'$, and applying
$\pi$ gives $a \Rightarrow_{\langle l, r \rangle} b$.

$(1) \Rightarrow (3)$.
Fix raw $A, L, R$ over $a, l, r$.
By~\cref{def:enriched_prop_rewriting} there are morphisms of
$\mathbf{H}^{\join}_{\Sigma}$ exhibiting $a$ and $b$ in a common context with $l$ and
$r$ in the hole; choosing raw representatives of the context and of $l, r$ gives raw
$A^{\circ} \Rightarrow_{\langle L^{\circ}, R^{\circ} \rangle} B^{\circ}$ with
$\pi(A^{\circ}) = a$, $\pi(B^{\circ}) = b$ and $\pi(L^{\circ}) = l$,
$\pi(R^{\circ}) = r$.
By~\cref{theorem:edpoi_raw},
$\mathcal{E}(\corner{A^{\circ}}) \;\overline{\rightsquigarrow}\;
\mathcal{E}(\corner{B^{\circ}})$ via
$\langle \mathcal{E}(\corner{L^{\circ}}), \mathcal{E}(\corner{R^{\circ}}) \rangle$.
Since $\pi(A) = \pi(A^{\circ})$, \cref{prop:interpretation_factorisation} gives
$\mathcal{E}(\corner{A}) \sim_{\mathfrak{R}} \mathcal{E}(\corner{A^{\circ}})$, and
likewise for $L, R$; replacing the rule sides by $\mathfrak{R}$-equivalent ones is
absorbed by~\cref{def:rewriting_modulo}.
Taking $B \defeq B^{\circ}$ yields (3).
\end{proof}

\begin{remark}
The asymmetry of clause~(3)---\emph{every} source representative, \emph{some} target
representative---is what makes the theorem usable for equality saturation: whatever
extended cospan currently represents $a$, the rule can be applied to it, and the result
represents $b$.
The modulo-$\mathfrak{R}$ is unavoidable at this level, because $\Rightarrow$ on
$\mathbf{H}^{\join}_{\Sigma}$ is by construction a relation on $\catname{SLat}$-classes;
by~\cref{remark:no_modulo_raw} it disappears entirely at the raw level, which is the
level at which one computes.
\end{remark}

\begin{remark}[Clause~(3) fails for $\overline{\rightsquigarrow}$]\label{remark:clause3_needs_modulo}
It is worth seeing that the modulo in clause~(3) cannot be dropped, even though it can be
dropped from~\cref{theorem:edpoi_raw}.
Take $a \defeq l$, so that $b = r$ and $[b] = [r]$, and consider the source
representative
\[
a' \defeq l \join l~, \qquad\text{for which}\qquad [a'] = [l] \join [l] = [l] = [a]
\]
by idempotency.
The only raw contexts exhibiting $l$ in $a'$ are $[-] \join l$ and $l \join [-]$, so a
single EDPOI step out of $\mathcal{E}(\corner{a'})$ produces
$\mathcal{E}(\corner{r \join l})$, and
\[
[r \join l] = [r] \join [l] \neq [r] = [b]
\]
in general---rewriting both components takes two steps.
With $\overline{\rightsquigarrow}^{\,\mathfrak{R}}$ the clause is restored, since
$\mathcal{E}(\corner{l \join l}) \sim_{\mathfrak{R}} \mathcal{E}(\corner{l})$ by the
idempotency rule of $\mathfrak{R}$, after which the rule applies directly.
The moral is that $\mathfrak{R}$ is what lets an arbitrary representative be brought into
a form exhibiting the redex; it is not an artefact of the formulation.
\end{remark}

\begin{remark}[On the purely existential form]\label{remark:existential_form}
Clause~(2) alone---four existentials, one for each corner---is equivalent to requiring
that \emph{some} extended cospan in the $\sim_{\mathfrak{R}}$-class of
$\mathcal{E}(\corner{a})$ rewrite to \emph{some} extended cospan in the class of
$\mathcal{E}(\corner{b})$, since by~\cref{theorem:concrete_representation}
and~\cref{prop:interpretation_factorisation}
\[
\bigl\{\, \mathcal{E}(\corner{a'}) \;\bigm|\; [a'] = [a] \,\bigr\}
\;=\;
\bigl\{\, \mathcal{X} \;\bigm|\; \mathcal{X} \sim_{\mathfrak{R}} \mathcal{E}(\corner{a}) \,\bigr\}~.
\]
It is therefore a faithful statement about $\catname{SLat}$-morphisms, but it says nothing
about any particular extended cospan, and in particular does not license applying a rule
to the e-hypergraph one happens to hold.
That is the content of clause~(3), and, in its sharpest form, of
\cref{theorem:edpoi_raw}.
\end{remark}

\begin{corollary}
\[
\mathbf{H}^{\join}_{\Sigma, \mathcal{E}}
\;\cong\;
\catname{EHypI}_{\bullet}(\Sigma) / \bigl(\mathfrak{R} \cup \mathfrak{R}_{\mathcal{E}}\bigr)~,
\qquad
\mathfrak{R}_{\mathcal{E}} \defeq
\bigl\{
\langle \mathcal{E}(\corner{L}), \mathcal{E}(\corner{R}) \rangle,
\langle \mathcal{E}(\corner{R}), \mathcal{E}(\corner{L}) \rangle
\;\bigm|\;
\pi(L) = l,\ \pi(R) = r,\ l = r \in \mathcal{E}
\bigr\}~.
\]
\end{corollary}
